\section{Related Work}
\label{sec:related_work}

\subsection{Layer-wise Readout and Internal State Analysis}

A growing body of work has developed tools for inspecting transformer internal states \citep[see][for a review]{bereska2024mechanistic}.
Logit Lens~\citep{nostalgebraist2020} and Tuned Lens~\citep{belrose2023eliciting} study how intermediate predictions emerge across layers; \citet{geva2021kv} and \citet{geva2022transformer} further show that feed-forward layers act as key-value memories that progressively promote concepts in the vocabulary space.
Patchscopes~\citep{ghandeharioun2024patchscopes} inject a source-layer representation into a target prompt and use the model's own generation process to decode what is encoded at that layer.
Linear probing~\citep{hewitt2019structural, hewitt2019designing, belinkov2022probing} tests whether specific attributes are linearly separable in hidden states, while recent work has extended representation-level analysis to broader properties such as spatial and temporal structure \citep{gurnee2024spacetime, burns2023discovering} and top-down control directions \citep{zou2023representation}.
Our work builds on this toolchain, but asks a different question: not only whether an answer is encoded at a layer, but whether it has already been organized into a form the model itself can use.
This is why we explicitly pair probing with CSD to distinguish information availability from information readiness.

\subsection{Understanding LLMs on Code}

Prior work on understanding how LLMs process code has largely focused on probing and attribution.
\citet{hooda2024counterfactual} use counterfactual analysis to test whether models truly respond to code predicates such as loop invariants and variable types, finding that models often rely on shallow pattern matching.
\citet{liu2024codemind} propose CodeMind, a framework that probes code reasoning across dependent and independent execution stages, revealing that LLMs struggle most when variable states depend on control-flow decisions.
\citet{gu2024cruxeval} introduce CRUXEval, a benchmark that evaluates code reasoning via input/output prediction, exposing significant gaps between code generation and code understanding capabilities.
\citet{chen2025reval} and \citet{abdollahi2025demystifying} further analyze how LLMs simulate code execution, documenting systematic errors on loops, arithmetic, and conditional branching---patterns consistent with our outcome fingerprints.
\textsc{AutoProbe}~\citep{autoprobe2025} dynamically aggregates informative layers and token positions to predict code correctness, showing that the most useful layers vary across models and tasks.
Neuron-Guided Interpretation of Code LLMs~\citep{neuronguided2026} moves toward neuron-level analysis in the code domain.
These works show that code-relevant information is not uniformly distributed inside models, but they mainly ask whether code properties are represented and where.
Our focus is instead on \emph{when} such information becomes self-decodable by the model, and how that timing changes systematically across code primitives.

\subsection{Reasoning Dynamics, Overthinking, and Causal Structure}

Our work is also related to recent studies of temporal reasoning dynamics and causal mechanism validation.
\citet{wendler2024llamas} reveal a phase structure in multilingual models, separating input encoding, conceptual processing, and output projection.
\citet{selfverify2025} show that models may internally encode correctness signals before these signals appear in output, and document token-level overthinking.
The overthinking phenomenon---where models degrade correct intermediate computations in later layers---has been characterized by \citet{halawi2024overthinking} and connected to early exit strategies \citep{schuster2022confident}; our Overprocessed category provides a fine-grained, task-specific instantiation of this phenomenon.
At the same time, \citet{wang2022ioi} and \citet{nikankin2025arithmetic} demonstrate how patching, ablation, and neuron-level analysis can provide causal evidence for internal mechanisms \citep{meng2022locating, geiger2023causal}, while \citet{universalneurons2024} and \citet{gould2024successor} study computational units shared across models.
\citet{nanda2023grokking} use mechanistic interpretability to track progress measures during training, showing that internal structure can be read off before behavioral changes emerge---an observation that parallels our finding that brewing precedes resolution.
We share the view that internal computation is structured and causally testable, but focus on layer-level reasoning dynamics in the code domain: answers often become externally readable before they become self-decodable, and both the brewing gap and the eventual resolution outcome depend strongly on the code primitive itself.
More broadly, \citet{nanda2025pragmatic} argue that code and tool-use are natural domains for interpretability because they provide verifiable ground truth; our setup follows this perspective by treating code execution primitives as the unit of analysis.
